\subsection{局部紧域上的有限维代数}

设 K 为交换域，A 为带单位元的有限秩 K-代数。对于每个 \(a \in A\)，令 \(L_{a}, R_{a}\) 分别为向量空间 A 的端同态 \(x \mapsto ax\)、\(x \mapsto xa\)，并令 \(\mathrm{N}_{\mathrm{A}/\mathrm{K}}(a) \in \mathrm{K}\)、\(\mathrm{N}_{\mathrm{A}^{0}/\mathrm{K}}(a) \in \mathrm{K}\) 为 a 在 A 的正则表示和相反代数 \(A^{0}\) 的正则表示中的范数；回顾 \(\mathrm{N}_{\mathrm{A}/\mathrm{K}}(a) = \det(L_{a})\)、\(\mathrm{N}_{\mathrm{A}^{0}/\mathrm{K}}(a) = \det(R_{a})\)（A, III, §9, No.~3）。以下条件等价：a 可逆，\(L_{a}\) 在 \(\operatorname{Hom}_{\mathrm{K}}(\mathrm{A}, \mathrm{A})\) 中可逆，\(R_{a}\) 在 \(\operatorname{Hom}_{\mathrm{K}}(\mathrm{A}, \mathrm{A})\) 中可逆，\(\mathrm{N}_{\mathrm{A}/\mathrm{K}}(a) \neq 0\)，\(\mathrm{N}_{\mathrm{A}^{0}/\mathrm{K}}(a) \neq 0\)。我们用 \(A^{*}\) 表示 A 的可逆元素集合。

现在设域 K 局部紧，从而代数 A 局部紧。于是 \(N_{A/K}\) 和 \(N_{A^{0}/K}\) 是从 A 到 K 的连续映射，故 \(A^{*}\) 是 A 中的开集。由第 10 号命题 15 的推论 1，

\[
\operatorname{mod} _ {\mathrm{A}} L _ {a} = \operatorname{mod} _ {\mathrm{K}} \mathrm{N} _ {\mathrm{A} / \mathrm{K}} (a), \quad \operatorname{mod} _ {\mathrm{A}} R _ {a} = \operatorname{mod} _ {\mathrm{K}} \mathrm{N} _ {\mathrm{A} ^ {0} / \mathrm{K}} (a).\tag{38}
\]

命题 16. --- 设 \(\alpha\) 为 A 的加法群上的哈尔测度。测度

\[
\left(\operatorname{mod} _ {K} N _ {A / K} (a)\right) ^ {- 1} d \alpha (a), \quad \left(\operatorname{mod} _ {K} N _ {A ^ {0} / K} (a)\right) ^ {- 1} d \alpha (a)
\]

在 \(\mathbf{A}^*\) 上分别是乘法群 \(\mathbf{A}^*\) 的左、右哈尔测度。

令 \(\alpha'\) 为 \(\alpha\) 在开集 \(\mathbf{A}^*\) 上的限制。对于 \(a \in \mathbf{A}^*\)，有 \(L_a(\alpha') = (\mathrm{mod}_K \mathrm{N}_{\mathrm{A}/\mathrm{K}}(a))^{-1} \alpha'\)，因此 \((\mathrm{mod}_K \mathrm{N}_{\mathrm{A}/\mathrm{K}}(a))^{-1} d\alpha'(a)\) 是 \(\mathbf{A}^*\) 上的左哈尔测度（第 8 号，命题 10 的推论 1）。转到相反代数，可见 \((\mathrm{mod}_K \mathrm{N}_{\mathrm{A}^0/\mathrm{K}}(a))^{-1} d\alpha'(a)\) 是 \(\mathbf{A}^*\) 上的右哈尔测度。

命题 17. --- 设 A 为域（局部紧）。对于每个 \(a \in \mathbf{A}\)，\(\operatorname{mod}_{\mathbf{A}}(a) = \operatorname{mod}_{\mathbf{K}} \mathrm{N}_{\mathbf{A}/\mathbf{K}}(a)\)。

这是（38）中第一个公式的翻译。

例。--- 1) 取 K = R，A = C。根据 Alg., chap.~VIII, §12, n° 2, 命题 4，得到 mod \(_{C}(z) = |z|^{2}\) 对所有 \(z \in \mathbf{C}\) 成立。(2)

\begin{enumerate}
\def\labelenumi{\arabic{enumi})}
\setcounter{enumi}{1}
\tightlist
\item
  取 \(\mathbf{K} = \mathbf{R}\)，令 \(\mathbf{A}\) 为四元数域 \(\mathbf{H}\)（GT, VIII, §1, No.~4）。考察 \(\mathbf{M}_2(\mathbf{C})\) 中以下元素：
\end{enumerate}

\[
X _ {1} = \left( \begin{array}{c c} 0 & i \\ i & 0 \end{array} \right) \quad X _ {2} = \left( \begin{array}{c c} 0 & - 1 \\ 1 & 0 \end{array} \right) \quad X _ {3} = \left( \begin{array}{c c} i & 0 \\ 0 & - i \end{array} \right)
\]

它们与 \(I_2\) 一起构成 \(\mathbf{M}_2(\mathbf{C})\) 上的 \(\mathbf{C}\) 的一组基。容易验证

\[
\begin{array}{l l} {X _ {1} ^ {2} = X _ {2} ^ {2} = X _ {3} ^ {2} = - I _ {2},} & {X _ {1} X _ {2} = - X _ {2} X _ {1} = X _ {3},} \\ {X _ {2} X _ {3} = - X _ {3} X _ {2} = X _ {1},} & {X _ {3} X _ {1} = - X _ {1} X _ {3} = X _ {2}.} \end{array}
\]

因此，映射 \(a + bi + cj + dk \mapsto aI_2 + bX_1 + cX_2 + dX_3\) 可延拓为代数 \(\mathbf{C} \otimes_{\mathbf{R}} \mathbf{H}\) 到代数 \(\mathbf{M}_2(\mathbf{C})\) 的 C-同构。由于 \([\mathbf{H} : \mathbf{R}] = 4\)，\(\mathbf{C}\) 是 \(\mathbf{H}\) 的分裂域（Alg., chap.~VIII, §10, n° 5），而 \(q = a + bi + cj + dk \in \mathbf{H}\) 的约化范数为

\[
\begin{array}{r l} \mathrm{Nrd} (q) & = \det (a I _ {2} + b X _ {1} + c X _ {2} + d X _ {3}) \\ & = \det \left( \begin{array}{c c} a + i d & - c + i b \\ c + i b & a - i d \end{array} \right) = a ^ {2} + b ^ {2} + c ^ {2} + d ^ {2} = \| q \| ^ {2}. \end{array}
\]

由 Alg., chap.~VIII, §12, n° 3, 命题 8，有

\[
\mathrm{N} _ {\mathbf {H} / \mathbf {R}} (q) = \left(\mathrm{Nrd} _ {\mathbf {H} / \mathbf {R}} (q)\right) ^ {2} = \| q \| ^ {4}.
\]

据此，命题 17 表明

\[
\operatorname{mod} _ {\mathbf {H}} (q) = \| q \| ^ {4}.
\]

关于局部紧域结构的更深入研究将在 CA, VI, §9 中进行。

